% =========================================================================
% English translation (generated by AI using Claude Opus 4.8 (Max effort)).
% This file was translated from Hungarian to English by AI.
% DISCLAIMER: Please review against the original Hungarian .tex files, before relying on it!
% SOURCE: 18. Számítógépes hálózatok és Internet eszközök.tex
% =========================================================================
\documentclass[margin=0px]{article}

\usepackage{listings}
\usepackage[utf8]{inputenc}
\usepackage{graphicx}
\usepackage{float}
\usepackage[a4paper, margin=0.7in]{geometry}
\usepackage{subcaption}
\usepackage{amsthm}
\usepackage{amssymb}
\usepackage{amsmath}
\usepackage{color}
\usepackage{fancyhdr}
\usepackage{setspace}

\onehalfspacing

\renewcommand{\figurename}{Figure}
\newenvironment{tetel}[1]{\paragraph{#1 \\}}{}
\newcommand{\R}{\mathbb{R}}

\pagestyle{fancy}
\lhead{\it{CS BSc Final Exam Topics}}
\rhead{18. Computer networks}

\title{\textbf{{\Large ELTE IK - Computer Science BSc} \vspace{0.2cm} \\ {\huge Final Exam Topics}} \vspace{0.3cm} \\ 18. Computer networks}
\author{}
\date{}

\begin{document}
\maketitle

\begin{center}
\fbox{\parbox{0.92\textwidth}{\small \textit{Note: This document was translated from Hungarian to English by AI. Please verify the information against the original Hungarian source before relying on it.}}}
\end{center}

\begin{tetel}{Computer networks}
    Layer models. Physical layer: baseband, broadband, digital encodings, modulation. Data link layer: framing, error control (detection, correction), CRC, flow control, dynamic channel allocation. Network layer: distance vector protocol, link-state protocol, BGP, path-vector protocol, IPv4 vs IPv6. Transport layer: UDP, TCP (connection management, congestion); Application layer: DNS, HTTP, DHCP, ARP.
\end{tetel}

\section{Models of networks}
\begin{description}
    \item[TCP/IP model]  \hfill \\
        \textbf{T}ransmission \textbf{C}ontrol \textbf{P}rotocol/\textbf{I}nternet \textbf{P}rotocol. Briefly TCP/IP. The TCP/IP model became the standard of American military-purpose computer networks in 1982. From 1985 it was popularized for commercial use.

        It distinguishes 4 layers:
        \begin{enumerate}
            \item Link layer
            \item Internet (network) layer
            \item Transport layer
            \item Application layer
        \end{enumerate}
    \item[OSI model]  \hfill \\
        \textbf{O}pen \textbf{S}ystem \textbf{I}nterconnection Reference Model. Briefly OSI reference model. It defines a standard conceptual model for the internal functionality of communication networks.

        It distinguishes 7 layers:
        \begin{enumerate}
            \item Physical layer
            \item Data link layer
            \item Network layer
            \item Transport layer
            \item Session layer
            \item Presentation layer
            \item Application layer
        \end{enumerate}
\end{description}
\section{Physical layer}
\begin{description}
    \item[Definition] \hfill \\
        The task of the physical layer is the transmission of bits over the communication channel. That is, ensuring correct bit transfer, the handling of the connection, and the clarification of the time and other details needed for the transmission. \\
        So the design considerations relate to the mechanical, electrical and procedural questions of the interface, and to the transmission medium.
    \item[Data transmission] \hfill
        \begin{description}
            \item[Wired] \hfill \\
                Data transmission in the case of a wire is possible by changing some physical characteristic (e.g.: voltage, current). This we can characterize with a periodic function $g(t)$.
                \begin{align*}
                    g(t) = \frac{1}{2}c+ \sum\limits_{n=1}^\infty a_nsin(2\pi n f t) + \sum\limits_{n=1}^\infty b_ncos(2\pi n f t)
                \end{align*}
                where $f=\frac{1}{T}$ is the fundamental frequency, and $a_n$ and $b_n$ are the $n$-th harmonic sine and cosine amplitudes
            \item[Wireless] \hfill \\
                For wireless data transmission, electromagnetic waves are used in many places. The waves have a frequency and a wavelength.
                \begin{itemize}
                    \item Frequency:
                          The number of oscillations of the wave per second. Notation: $f$, unit: Hz (Hertz)
                    \item Wavelength: the distance between two consecutive wave crests (or wave troughs). Notation: $\lambda$
                \end{itemize}
                \begin{align*}
                    \lambda f = c
                \end{align*}
                where $c$ is the speed of light, that is, the propagation speed of electromagnetic waves in vacuum.
                \begin{figure}[H]
                    \centering
                    \includegraphics[width=0.6\textwidth]{../../18. Számítógépes hálózatok és Internet eszközök/img/elektromagneses.png}
                    \caption{Electromagnetic spectrum}
                \end{figure}
            \item[Symbols] \hfill \\
                Instead of bits we transmit symbols. (E.g. 4 symbols: A - 00, B - 01, C - 10, D - 11)\\
                Baud: symbol/second \\
                Data rate: bit/second
            \item[Synchronization] \hfill \\
                Question: When do signals have to be measured, and when does a symbol begin? For this synchronization is needed between the parties.
                \begin{itemize}
                    \item Explicit clock \\
                          Use of parallel transmission channels, synchronized data, suitable in case of short transmission.
                    \item Critical time points \\
                          Let us synchronize, for example, at the start of a symbol or block; outside the critical time points the clocks run freely, we assume that the clocks run in sync for a short time.
                    \item Symbol codes \\
                          Self-clocking signal -- a signal decodable without separate clock synchronization, the signal contains the information needed for synchronization.
                \end{itemize}
                \begin{figure}[H]
                    \centering
                    \includegraphics[width=0.6\textwidth]{../../18. Számítógépes hálózatok és Internet eszközök/img/szinkronizacio.png}
                    \caption{The necessity of synchronization}
                \end{figure}
        \end{description}
    \item[Transmission media] \hfill \\
        \begin{description}
            \item[Wired] \hfill
                \begin{itemize}
                    \item magnetic data carriers \\
                          bandwidth good, latency large
                    \item twisted pair \\
                          Mainly used in telephone systems; double copper wire; analog and digital signal transmission
                    \item Coaxial cable\\
                          Higher speed and distance can be achieved than with twisted pair; analog and
                          digital signal transmission
                    \item Optical fibers\\
                          Light source, transmission medium and detector; light pulse a 1 bit, no
                          light pulse a 0 bit
                \end{itemize}
            \item[Wireless] \hfill
                \begin{itemize}
                    \item Radio frequency transmission \\
                          easily producible; large distance; outdoor and indoor applicability; frequency-dependent propagation characteristics
                    \item Microwave transmission\\
                          propagates along a straight line; problem of fading; not expensive
                    \item Infrared and millimeter-wave transmission \\
                          for short-distance transmission; does not penetrate solid objects
                    \item Visible light wave transmission\\
                          laser source + light sensor; large bandwidth, cheap, license-free; weather can strongly affect it
                \end{itemize}
        \end{description}
    \item[Signal transmission] \hfill
        \begin{itemize}
            \item Baseband \\
                  The digital signal is directly transformed into current or voltage. The signal is transmitted on every frequency. Transmission limits
                  \begin{figure}[H]
                      \centering
                      \includegraphics[width=0.8\textwidth]{../../18. Számítógépes hálózatok és Internet eszközök/img/alapsav.png}
                      \caption{Structure of digital baseband transmission}
                  \end{figure}
            \item Broadband \\
                  The transmission happens in a wide frequency range. For modulating the signal we can use the following possibilities.
                  \begin{figure}[H]
                      \centering
                      \includegraphics[width=0.8\textwidth]{../../18. Számítógépes hálózatok és Internet eszközök/img/szelessav.png}
                      \caption{Structure of digital broadband transmission}
                  \end{figure}

                  Modulations:\\
                  The representation of a sinusoidal oscillation as a function of period $T$ is $g(t)=Asin(2\pi f t + \varphi)$, where $A$ is the amplitude, $f = \frac{1}{T}$ the frequency and $\varphi$ the phase shift
                  \begin{itemize}
                      \item Amplitude modulation \\
                            We encode the signal $s(t)$ as the amplitude of the sine curve, that is:
                            \begin{align*}
                                f_A(t) = s(t) \cdot sin(2\pi f t + \varphi)
                            \end{align*}
                            Analog signal: amplitude modulation\\
                            Digital signal: amplitude keying (the strength of the signal changes according to the values of a discrete set)
                            \begin{figure}[H]
                                \centering
                                \includegraphics[width=0.5\textwidth]{../../18. Számítógépes hálózatok és Internet eszközök/img/amplitudo_mod.png}
                                \caption{Amplitude modulation}
                            \end{figure}
                            \begin{figure}[H]
                                \centering
                                \includegraphics[width=0.5\textwidth]{../../18. Számítógépes hálózatok és Internet eszközök/img/amplitudo_key.png}
                                \caption{Amplitude keying}
                            \end{figure}
                      \item Frequency modulation \\
                            We encode the signal $s(t)$ in the frequency of the sine curve, that is:
                            \begin{align*}
                                f_F(t) = a \cdot sin(2\pi s(t) t + \varphi)
                            \end{align*}
                            Analog signal: frequency modulation\\
                            Digital signal: frequency-shift keying (for example by assigning different frequencies to the symbols of a discrete set)
                            \begin{figure}[H]
                                \centering
                                \includegraphics[width=0.5\textwidth]{../../18. Számítógépes hálózatok és Internet eszközök/img/frekvencia_mod.png}
                                \caption{Frequency modulation}
                            \end{figure}
                            \begin{figure}[H]
                                \centering
                                \includegraphics[width=0.5\textwidth]{../../18. Számítógépes hálózatok és Internet eszközök/img/frekvencia_key.png}
                                \caption{Frequency keying}
                            \end{figure}
                      \item Phase modulation \\
                            We encode the signal $s(t)$ in the phase of the sine curve, that is:
                            \begin{align*}
                                f_P(t) = a \cdot sin(2\pi f t + s(t))
                            \end{align*}
                            Analog signal: phase modulation (not really used)\\
                            Digital signal: phase-shift keying (for example by assigning different phases to the symbols of a discrete set)
                            \begin{figure}[H]
                                \centering
                                \includegraphics[width=0.5\textwidth]{../../18. Számítógépes hálózatok és Internet eszközök/img/fazis_mod.png}
                                \caption{Phase modulation}
                            \end{figure}
                            \begin{figure}[H]
                                \centering
                                \includegraphics[width=0.5\textwidth]{../../18. Számítógépes hálózatok és Internet eszközök/img/fazis_key.png}
                                \caption{Phase-shift keying}
                            \end{figure}
                  \end{itemize}

                  Comparison of digital and analog signals: \\
                  Digital transmission: Uses a finite set of discrete signals
                  (for example voltage or current
                  values).\\
                  Analog transmission: Uses a continuous set of signals
                  (for example voltage or current in the wire) \\
                  In the case of digital there is the possibility of restoring the reception accuracy and restoring the original signal, while for analog the arising errors can amplify themselves.
        \end{itemize}
\end{description}
\section{Data link layer}
\begin{description}
    \item[Definition] \hfill \\
        The task of the data link layer is to provide a well-defined service interface to the network layer, which has three phases:
        \begin{itemize}
            \item unacknowledged connection-based service
            \item acknowledged connectionless service
            \item acknowledged connection-based service
        \end{itemize}
        Furthermore the handling of transmission errors and the regulation of the data traffic (avoidance of flooding).
    \item[Framing] \hfill \\
        The physical layer does not guarantee error-freeness; the task of the data link layer is error signaling and, as needed, correction. The solution for this: breaking the bit stream into frames, and computing checksums. Framing is not a simple task, since there is not much possibility for reliable timing. Four possible methods:
        \begin{enumerate}
            \item Character counting \\
                  We give the number of characters in the frame in the header of the frame. Thus the data link layer of the receiver will know the end of the frame. Problem: the method is very sensitive to error.
                  \begin{figure}[H]
                      \centering
                      \includegraphics[width=0.7\textwidth]{../../18. Számítógépes hálózatok és Internet eszközök/img/karakterszamlalas.png}
                      \caption{Character counting}
                  \end{figure}
            \item Start and end characters with character stuffing \\
                  We place special bytes to indicate the start and end of the frame, whose name is flag byte. For the special bytes appearing in the data stream an ESC byte is used.
                  \begin{figure}[H]
                      \centering
                      \includegraphics[width=0.6\textwidth]{../../18. Számítógépes hálózatok és Internet eszközök/img/karakterbeszuras.png}
                      \caption{Start and end characters with character stuffing}
                  \end{figure}
            \item Start and end marks with bit stuffing \\
                  Every frame starts with a special bit pattern (flag byte, 01111110) and after every consecutive continuous run of 5 1-bits it inserts a 0.
                  \begin{figure}[H]
                      \centering
                      \includegraphics[width=0.9\textwidth]{../../18. Számítógépes hálózatok és Internet eszközök/img/bitbeszuras.png}
                      \caption{Start and end marks with bit stuffing}
                  \end{figure}
            \item Physical-layer coding violation \\
                  Usable in networks where the physical-layer coding contains redundancy.
        \end{enumerate}
    \item[Error handling] \hfill \\
        From the point of view of error handling we have to examine the following two cases. Whether the frames arrived at the network layer of the destination, and whether they arrived in the correct order. For this some feedback is needed between the receiver and the sender. (for example acknowledgments). We introduce time limits for the individual steps. In case of an error we resend the packet. Multiple reception can occur, with which the use of sequence numbers can help. The task of the data link layer, from the point of view of error handling, is to handle the timers and counters so that it can ensure the exactly-once (not more and not less) arrival of the frames at the network layer of the destination.

        Bit errors:
        \begin{itemize}
            \item single bit error \\
                  1 bit of the data unit changes from zero to one or from one to zero.
                  \begin{figure}[H]
                      \centering
                      \includegraphics[width=0.2\textwidth]{../../18. Számítógépes hálózatok és Internet eszközök/img/egyszeru_bithiba.png}
                      \caption{Single bit error}
                  \end{figure}
            \item burst bit error \\
                  A continuous symbol sequence whose first and last symbols are erroneous, and in the subsequence bounded by these two symbols there is no subsequence of length $m$ that we received correctly, is called a burst bit error of length $m$.
            \end{itemize}

        Error detection and correction:\\
        There are two kinds of error-handling strategy. These are the error-detecting (attaching redundant information) and error-correcting codes (redundancy inserted among the data).
            [On reliable channels error detection is cheaper. (we rather resend the packet). On less reliable channels the error-correcting method is more practical]
        \begin{itemize}
            \item Hamming distance, Hamming bound \\
                  The frame to be sent contains $m$ bits. The number of redundant bits is $r$. So the sent frame is: $n = m+r$ bits.

                  Hamming distance: \\
                  The number of bit positions on which different bits stand in the two code words is called the Hamming distance of the two code words. Notation: $d(x,y)$\\
                  Let $S$ be a set of equal-length bit words. The Hamming distance of $S$:
                  \begin{align*}
                      d(S) := \min_{x,y, \in S \ \land \ x \neq y} d(x,y)
                  \end{align*}

                  In case d(S) = 1: \\
                  There is no error detection, since from an allowed code word, by changing 1 bit, an allowed code word is produced.
                  \begin{figure}[H]
                      \centering
                      \includegraphics[width=0.4\textwidth]{../../18. Számítógépes hálózatok és Internet eszközök/img/hamming1.png}
                      \caption{Hamming distance of the code = 1}
                  \end{figure}
                  In case d(S) = 2: \\
                  If for the code word $x$ there exists a non-allowed code word $v$ for which $d(u,x)=1$, then an error occurred. If $x$ and $y$ are allowed code words (their distance is minimal = 2), then the following relationship has to hold:
                  \begin{align*}
                      2 = d(x,y) \leq d(x,v)+d(v,y)
                  \end{align*}
                  That is, a bit error is detectable but not correctable.
                  \begin{figure}[H]
                      \centering
                      \includegraphics[width=0.4\textwidth]{../../18. Számítógépes hálózatok és Internet eszközök/img/hamming2.png}
                      \caption{Hamming distance of the code = 2}
                  \end{figure}

                  In case d(S) = 3: \\
                  Then every $u$ for which $d(x,u)=1$ and $d(u,y) > 1$ is not allowed. Then three possibilities hold:
                  \begin{itemize}
                      \item x was transmitted and arrived with 1 bit error
                      \item y was transmitted and arrived with 2 bit errors
                      \item something else was transmitted and arrived with at least 2 bit errors
                  \end{itemize}
                  But it is more probable that $x$ was transmitted, so this is a 1-bit-error-correcting, 2-bit-error-detecting code.
                  \begin{figure}[H]
                      \centering
                      \includegraphics[width=0.4\textwidth]{../../18. Számítógépes hálózatok és Internet eszközök/img/hamming3.png}
                      \caption{Hamming distance of the code = 3}
                  \end{figure}

                  Hamming bound: \\
                  $ C \subseteq \{0,1\}^n$ and $d(C) = k$. Then the union of the pairwise disjoint sets of bit words found in the radius-$\frac{k-1}{2}$ neighborhoods of the code words can produce at most the set of $n$-length bit words. That is, formally:
                  \begin{align*}
                      |C|\sum_{i=0}^{\lfloor\frac{k-1}{2}\rfloor}\dbinom{n}{i} \leq 2^n
                  \end{align*}

                  \begin{itemize}
                      \item Error detection: \\
                            For the detection of $d$ bit errors, at least $d+1$ Hamming distance is needed in the set of frames.

                      \item Error correction: \\
                            For the correction of $d$ bit errors, at least $2d+1$ Hamming distance is needed in the set of allowed frames.

                      \item Rate of the code: \\
                            $R_S = \frac{log_2|S|}{n}$ is the rate of the code ($S \subseteq \{0,1\}^n$) - characterizes efficiency

                      \item Distance of the code: \\
                            $\delta_S = \frac{d(S)}{n}$ is the distance of the code ($S \subseteq \{0,1\}^n$) - characterizes error handling
                  \end{itemize}
                  A good code has both a large rate and a large distance.
            \item Parity bit \\
                  The parity bit is a bit that we choose based on the number of ones in the code word.
                  \begin{itemize}
                      \item Odd Parity - if the number of ones is odd, then inserting a 0; otherwise inserting a 1
                      \item Even Parity - if the number of ones is even, then inserting a 0; otherwise inserting a 1
                  \end{itemize}
                  A method using parity is the so-called Hamming method:\\
                  Let us number the bits of the code word (starting with 1). On the power-of-2 positions the check bits get a place, into the remaining places the message bits go. Every check bit sets the parity of a group of bits (including itself) to even (or odd).\\
                  The groups are formed as follows:
                  \begin{itemize}
                      \item 1st bit: Every first one-long bit sequence starting from the first bit (so: 1,3,5,7,...)
                      \item 2nd bit: Every first two-long bit sequence starting from the second bit (so: 2-3,6-7,10-11)
                      \item 4th bit: Every first four-long bit sequence starting from the fourth bit (so: 4-7,12-15)
                      \item etc.
                  \end{itemize}
                  \begin{figure}[H]
                      \centering
                      \includegraphics[width=0.8\textwidth]{../../18. Számítógépes hálózatok és Internet eszközök/img/paritasbit.png}
                      \caption{Groups of parity bits}
                  \end{figure}
                  Example:\\
                  Let the message to be sent be: 1000101 \\
                  Then the code word is formed as follows: \\
                  {\color[rgb]{1,0,0}$\heartsuit\heartsuit$}1{\color[rgb]{1,0,0}$\heartsuit$}000{\color[rgb]{1,0,0}$\heartsuit$}101\\
                  The 8th bit sets the parity of the 8-11 bit sequence to even:\\
                  {\color[rgb]{1,0,0}$\heartsuit\heartsuit$}1{\color[rgb]{1,0,0}$\heartsuit$}000{\color[rgb]{0,0,1}0101}\\
                  The 4th bit sets the parity of the 4-7 bit sequence to even:\\
                  {\color[rgb]{1,0,0}$\heartsuit\heartsuit$}1{\color[rgb]{0,0,1}0000}0101\\
                  The 2nd bit sets the parity of the 2-3, 6-7, 10-11 bit sequence to even:\\
                  {\color[rgb]{1,0,0}$\heartsuit$}{\color[rgb]{0,0,1}01}00{\color[rgb]{0,0,1}00}01{\color[rgb]{0,0,1}01}\\
                  The 1st bit sets the parity of the 1,3,5,7,9,11 bit sequence to even: \\
                  {\color[rgb]{0,0,1}1}0{\color[rgb]{0,0,1}1}0{\color[rgb]{0,0,1}0}0{\color[rgb]{0,0,1}0}0{\color[rgb]{0,0,1}1}0{\color[rgb]{0,0,1}1}\\

                  So the bit sequence to be sent:
                  {\color[rgb]{1,0,0}10}1{\color[rgb]{1,0,0}0}000{\color[rgb]{1,0,0}0}101\\

            \item CRC - Polynomial code, that is, cyclic redundancy \\
                  We regard the bit sequences as the coefficients of a polynomial ($M(x)$) over $\mathbb{Z}_2$. We define a generator polynomial $G(x)$ of degree $r$, which the receiver and sender both know.
                  \begin{enumerate}
                      \item Let us append $r$ 0 bits to the low-place-value end of the frame. That is, let us take the polynomial $x^rM(x)$ (this is already of degree $m+r$)
                      \item Let us divide $x^rM(x)$ by $G(x)$ (mod 2).
                      \item Let us subtract the remainder (which always contains r or fewer bits) from $x^rM(x)$ (mod 2). Thus an r-long checksum gets to the end of the original frame. Let this polynomial be $T(x)$.
                      \item The receiver gets a polynomial corresponding to $T(x)+E(x)$ (where $E(x)$ is the error polynomial). Dividing this by the generator polynomial we get a polynomial $R(x)$. If this polynomial is not zero, then an error occurred.
                  \end{enumerate}
                  The bit errors corresponding to multiples of $G(x)$ we do not detect.
        \end{itemize}
    \item[Protocols] \hfill
        \begin{description}
            \item[Elementary data link protocols] \hfill
                \begin{itemize}
                    \item Unrestricted simplex protocol\\
                          Environment:
                          \begin{itemize}
                              \item Sender, receiver: always ready.
                              \item No processing time
                              \item Infinite buffer
                              \item No frame corruption, loss
                          \end{itemize}

                          Protocol:
                          \begin{itemize}
                              \item No sequence number/acknowledgment
                              \item The sender sends the frames out continuously in an infinite loop
                              \item The receiver initially waits for the arrival of the first frame, at the arrival of a frame it puts the content of the hardware buffer into a variable and forwards the data part to the network layer.
                          \end{itemize}
                    \item Simplex stop-and-wait protocol 	\\
                          Environment:
                          \begin{itemize}
                              \item Network layers of sender, receiver: always ready.
                              \item The receiver needs $\delta t$ time for processing the incoming frame.
                              \item No buffering and no queuing either
                              \item No frame corruption, loss
                          \end{itemize}

                          Protocol:
                          \begin{itemize}
                              \item No sequence number/acknowledgment
                              \item The sender sends one by one, the next only following the acknowledgment.
                              \item The receiver initially waits for the arrival of the first frame, at the arrival of a frame it puts the content of the hardware buffer into a variable and forwards the data part to the network layer, finally it acknowledges the frame.
                          \end{itemize}
                    \item Simplex protocol for a noisy channel\\
                          Environment:
                          \begin{itemize}
                              \item Network layers of sender, receiver: always ready.
                              \item The receiver needs $\delta t$ time for processing the incoming frame.
                              \item No buffering and no queuing either
                              \item Frame can be corrupted, lost
                          \end{itemize}

                          Protocol:
                          \begin{itemize}
                              \item No sequence number/acknowledgment
                              \item The sender sends one by one, it does not send a new one until it gets an acknowledgment within the deadline. After the deadline it resends the frame.
                              \item The receiver initially waits for the arrival of the first frame, at the arrival of a frame it puts the content of the hardware buffer into a variable, checks the control sum:
                                    \begin{itemize}
                                        \item No error: it forwards the data part to the network layer, finally acknowledges the frame.
                                        \item There is an error: it discards the frame and does not acknowledge
                                    \end{itemize}
                          \end{itemize}
                \end{itemize}
            \item[Sliding window protocol] \hfill \\
                At a given moment several frames can be in transmission state at once. The receiver allocates a buffer of size corresponding to $n$ frames. The sender is allowed to send at most $n$, that is, window-size, unacknowledged frames. The position of the frame in the sequence gives the label of the frame (sequence number).
                The receiver has to discard the erroneous frames, and the frames arriving with a non-allowed sequence number. The sender keeps track of the set of sendable sequence numbers (sending window). The receiver keeps track of the set of receivable sequence numbers (receiving window). The sending window narrows with every sending, and grows with the arrival of an acknowledgment.

                What if a frame error occurs in the middle of a long stream?
                \begin{enumerate}
                    \item ``go-back-N'' strategy\\
                          It discards every frame after the erroneous frame and does not send an acknowledgment about them. When the sender's timer expires, it resends all unacknowledged frames, starting with the corrupted or lost frame. Disadvantage: It can waste a lot of bandwidth if the error rate is high.
                    \item ``selective repeat'' strategy \\
                          It discards the erroneous frames, but buffers the good frames after the erroneous one. When the sender's timer expires, it resends the oldest unacknowledged frame. Disadvantage: Large memory requirement in case of a large receiving window.
                \end{enumerate}

            \item[Examples of data link protocols] \hfill
                \begin{itemize}
                    \item HDLC - High-level Data Link Control \\
                          The HDLC protocol uses a 3-bit sliding-window protocol for the sequence numbering.
                          It uses three types of frame:
                          \begin{itemize}
                              \item information
                              \item supervisory
                                    \begin{itemize}
                                        \item acknowledgment frame (RECEIVE READY)
                                        \item negative acknowledgment frame (REJECT)
                                        \item not ready to receive (RECEIVE NOT READY) - acknowledges every frame up to the next
                                        \item selective reject (SELECTIVE REJECT) - calls for the resending of a given frame
                                    \end{itemize}
                              \item Unnumbered
                          \end{itemize}

                          General frame structure:
                          \begin{itemize}
                              \item FLAG byte for indicating the frame boundaries
                              \item \textit{address} field - has significance for terminals with several lines
                              \item \textit{control} field - for the tasks of sequence numbering, acknowledgment and others
                              \item \textit{data} field - can be data of arbitrary length
                              \item \textit{checksum} field - CRC control sum (using the CRC-CCITT generator polynomial)
                          \end{itemize}
                          \begin{figure}[H]
                              \centering
                              \includegraphics[width=0.8\textwidth]{../../18. Számítógépes hálózatok és Internet eszközök/img/hdlc_keret.png}
                              \caption{Structure of the HDLC frame}
                          \end{figure}
                    \item PPP - Point-to-Point Protocol \\
                          The PPP protocol ensures three things:
                          \begin{itemize}
                              \item  A framing method (unambiguous frame boundaries)
                              \item A link control protocol (for waking up, testing the lines, negotiating the options and releasing the lines.)
                              \item A way of negotiating the network-layer options that is independent of the applied network-layer protocol.
                          \end{itemize}
                          It uses a byte-based frame structure (that is, the smallest data unit is the byte).
                          The default value of the control field indicates the unnumbered frame. In the protocol field the protocol code can be for LCP, NCP, IP, IPX, AppleTalk or another protocol.
                          \begin{figure}[H]
                              \centering
                              \includegraphics[width=0.8\textwidth]{../../18. Számítógépes hálózatok és Internet eszközök/img/ppp_keret.png}
                              \caption{Structure of the PPP frame}
                          \end{figure}
                \end{itemize}
            \item[MAC - Media Access Control] \hfill \\
                In our discussions so far we assumed a point-to-point connection. Now we will deal with the subject of networks using a broadcast channel. The channel allocation can happen in a static or dynamic way.

                In the static case either Frequency Division Multiplexing or Time Division Multiplexing is used. In the frequency-division case the bandwidth is divided into $N$ parts, and every user gets one band. In the time-division case the time unit is divided into $N$ parts and these are given to the users. Both methods cannot be efficient in case of bursty traffic.

                In what follows we examine the dynamic channel allocation methods.
                \begin{itemize}
                    \item Contention protocols \\
                          There are N independent stations on which a program or a user generates frames to be transmitted. If a station generated a frame, then it stays in a blocked state until it successfully transmits the frame. There is a single channel on which all kinds of communication happen. Every station can send and receive data on this channel. If two frames are transmitted at the same time, then they overlap, and the resulting signal becomes uninterpretable. This is called a collision. This is detectable for every station. The frames involved in the collision have to be resent later. (Apart from this error no other error can occur.)

                          We distinguish two kinds of time model:
                          \begin{enumerate}
                              \item Continuous -- Every station can begin the transmission of its ready-to-send frame at an arbitrary time point.
                              \item Discrete -- We divide the time into discrete slots. Frame transmission is possible only at the start of a slot. The slot can be empty, successful or collided
                          \end{enumerate}

                          The individual stations either have carrier sensing or not. If not, then the stations cannot examine the state of the common channel, so they simply begin to send if they have the opportunity. If yes, then the stations can examine the state of the common channel before sending. The channel can be: busy or free. If the channel is busy, then the stations do not try to use it until it is freed

                          \begin{itemize}
                              \item Pure ALOHA \\
                                    The user can transfer data whenever they want. In case of a collision the station waits for a random time, then tries again.

                                    Frame time -- the time needed to transmit a standard, fixed-length frame

                                    A frame does not suffer a collision if, from the first moment of its sending, for one frame time no other station attempts a frame sending.

                                    \begin{figure}[H]
                                        \centering
                                        \includegraphics[width=0.6\textwidth]{../../18. Számítógépes hálózatok és Internet eszközök/img/egyszeru_aloha.png}
                                        \caption{Pure ALOHA frame collisions}
                                    \end{figure}
                              \item Slotted ALOHA \\
                                    By dividing the time into discrete slots aligned to the frame time, the capacity of the ALOHA system can be doubled. Even a small increase of the channel's load can drastically decrease the performance of the medium.
                              \item 1-persistent CSMA \\
                                    There is carrier sensing, that is, every station can listen in on the channel. The protocol uses a continuous time model.

                                    Algorithm:
                                    \begin{enumerate}
                                        \item Before transmitting a frame it listens in on the channel:
                                              \begin{enumerate}
                                                  \item If busy, then it waits until it is freed. In case of a free channel it sends immediately. (persistent)
                                                  \item If free, then it sends.
                                              \end{enumerate}
                                        \item If a collision occurs, then the station waits for a random time, then restarts the transmission of the frame.
                                    \end{enumerate}
                              \item Non-persistent CSMA \\
                                    There is carrier sensing, that is, every station can listen in on the channel. The protocol uses a continuous time model. It avoids greediness, that is, it does not send immediately if busy.

                                    Algorithm:
                                    \begin{enumerate}
                                        \item Before transmitting a frame it listens in on the channel:
                                              \begin{enumerate}
                                                  \item If busy, then it waits for a random time (does not watch the traffic), then restarts the sending algorithm. (non-persistent)
                                                  \item If free, then it sends.
                                              \end{enumerate}
                                        \item If a collision occurs, then the station waits for a random time, then restarts the transmission of the frame.
                                    \end{enumerate}
                              \item p-persistent CSMA \\
                                    There is carrier sensing, that is, every station can listen in on the channel. The protocol uses a discrete time model.

                                    Algorithm:
                                    \begin{enumerate}
                                        \item In ready-to-send state the station listens in on the channel:
                                              \begin{enumerate}
                                                  \item If busy, then it waits for the next slot, then repeats the algorithm.
                                                  \item If free, then with probability $p$ it sends, and with probability $1-p$ it backs off from its intention until the next slot. In case of waiting it repeats the algorithm in the next slot. This continues until it sends the frame, or until another station begins to send, because in this case it behaves as if a collision had occurred.
                                              \end{enumerate}
                                        \item If a collision occurs, then the station waits for a random time, then restarts the transmission of the frame.
                                    \end{enumerate}
                              \item CSMA/CD \\
                                    In case of collision detection it should be possible to interrupt the transmission. Every station observes the channel while sending, if it noticed a collision, then it interrupts the transmission, and waits for a random time, then begins to transmit its frame again
                          \end{itemize}
                          \begin{figure}[H]
                              \centering
                              \includegraphics[width=0.7\textwidth]{../../18. Számítógépes hálózatok és Internet eszközök/img/aloha_csma.jpg}
                              \caption{Comparison of the ALOHA and CSMA protocols}
                          \end{figure}
                    \item Collision-free protocols \\
                          Motivation: The collisions adversely affect the performance of the system, and CSMA/CD is not applicable everywhere.

                          There are N stations. The stations are unambiguously numbered from 0 to N. We assume a slotted time model.
                          \begin{itemize}
                              \item A reservation protocol \\
                                    If the i-th station wants to send, then in the i-th contention slot it can indicate this by sending a 1 bit. Thus by the end of the contention period every station knows the senders. The sending happens in the order of the sequence numbers.
                              \item Binary countdown protocol \\
                                    Every station has a binary identifier of equal length. The station wishing to transmit begins to send its binary address bit by bit, starting with the highest-place-value bit. The bits of the same position enter into a logical OR connection in case of a collision. If the station sends a zero but hears back a one, then it gives up its sending intention, because there is a sender with a larger identifier than it.
                          \end{itemize}
                    \item Limited contention protocols \\
                          A protocol that, in case of small load, uses a contention technique for the sake of small latency, and under large load applies a collision-free technique for the sake of good channel utilization.
                          \begin{itemize}
                              \item Adaptive tree walk \\
                                    Every station is represented by an unambiguous binary ID. The IDs correspond to the leaves of a (binary) tree. The slots are assigned to the individual nodes of the tree. In every slot we examine the subtree under the given node. At a node $u$ of the tree we can distinguish 3 cases:
                                    \begin{itemize}
                                        \item No station sends in the subtree $u$.
                                        \item Exactly one station sends in the subtree $u$.
                                        \item Several stations send in the subtree $u$. This is called a collision.
                                    \end{itemize}
                                    In case of a collision let us perform the check on both the left and right child of $u$.

                                    With this method it is easy to determine which station can send in the given slot.
                                    \begin{figure}[H]
                                        \centering
                                        \includegraphics[width=0.5\textwidth]{../../18. Számítógépes hálózatok és Internet eszközök/img/adaptivfa.png}
                                        \caption{Binary tree of the adaptive tree walk protocol}
                                    \end{figure}
                          \end{itemize}
                \end{itemize}
        \end{description}
\end{description}
\section{Network layer}
\begin{description}
    \item[Definition] \hfill \\
        The main task of the network layer is the forwarding of packets between the source and the destination. This is the lowest layer that deals with the transfer between two endpoints. It has to know the topology of the communication subnetwork. Care has to be taken not to overload certain communication paths, nor certain routers, in such a way that others stay idle.

        The services provided toward the transport layer:
        \begin{itemize}
            \item Independent of the construction of the subnetworks
            \item Hide the number, type and topology of the present subnetworks
            \item The network addresses made available to the transport layer have to form a uniform numbering system
        \end{itemize}
    \item[Types of routing] \hfill
        \begin{itemize}
            \item Hierarchical routing\\
                  We divide the routers into domains. Every router knows its own domain, but has no knowledge of the internal structure of the others. For large networks a multi-level hierarchy can be needed.
            \item Broadcast routing\\
                  the simultaneous sending of a packet everywhere.
            \item Multicast routing\\
                  The simultaneous sending of a packet to a specified group.

        \end{itemize}

    \item[Routing algorithms] \hfill \\
        That part of the network layer software that is responsible for the decision of on which output line the incoming packet gets forwarded. The process can be broken into two steps:
        \begin{enumerate}
            \item Filling and maintaining routing tables.
            \item Forwarding
        \end{enumerate}

        The classes of routing algorithms:
        \begin{enumerate}
            \item Adaptive algorithms
                  \begin{enumerate}
                      \item distance-based
                      \item link-based
                  \end{enumerate}
                  The topology and usually also the traffic can influence the decision.
            \item Non-adaptive algorithms \\
                  Offline determination, loading into the routers at start
        \end{enumerate}

        \begin{description}
            \item[Dijkstra algorithm] \hfill \\
                The Dijkstra algorithm is a static algorithm whose goal is the determination of the shortest path between two nodes.

                We label every node with the length of the shortest path found so far from the start point. During the operation of the algorithm the labels can change as paths are found. We distinguish two kinds of label: temporary and permanent. Initially every label is temporary. At the finding of the shortest path the label becomes a permanent label, and no longer changes.
            \item[Flooding algorithm] \hfill \\
                The flooding algorithm is a static algorithm.

                We forward every incoming packet on every outgoing line except the one on which it arrived. This way, however, very many duplicates would arise. Therefore
                \begin{itemize}
                    \item We introduce a hop counter, which every station decreases by one. If it decreases to 0, they discard it.
                    \item The stations keep track of the already-sent packets. Thus they do not send a packet several times.
                \end{itemize}
            \item[Distributed Bellman-Ford algorithm] \hfill \\
                The Distributed Bellman-Ford algorithm is an adaptive, distance-based routing algorithm. Every node can communicate only with its direct neighbors. Every station has its own distance vector. This it periodically sends to the direct neighbors. Every router knows the cost to its direct neighbors. Based on the received distance vectors every node updates its own vector.
            \item[Link-state routing] \hfill \\
                The motivation of the link-state routing algorithm is that the distance-based algorithms converged slowly, and the taking into account of differing bandwidths. \\
                The steps of the link-state routing algorithm:
                \begin{enumerate}
                    \item Discovering the neighbors, and determining their network addresses
                    \item Measuring the latency or cost to every neighbor
                    \item Assembling a packet from the learned information
                    \item Sending the packet to all the other routers
                    \item Computing the shortest path to all the other routers.
                \end{enumerate}
        \end{description}
    \item[Network layer on the Internet] \hfill \\
        At the level of the network layer the internet can be regarded as a connected collection of autonomous systems. It has no real structure, but numerous major backbone networks exist. To the backbone networks the territorial and regional networks connect. To the regional and territorial networks connect the LANs at universities, companies and internet providers.

        The protocol of the internet, the IP.

        On the Internet the communication works in the following way:
        \begin{enumerate}
            \item The transport layer carries the data streams and breaks them into datagrams.
            \item Every datagram gets transmitted on the Internet, possibly broken into smaller units along the way.
            \item The network layer of the destination machine assembles the original datagram, then passes it to its transport layer.
            \item The transport layer of the destination machine inserts the datagram into the input data stream of the receiving process.
        \end{enumerate}

        Internet Protocol - IP
        \begin{itemize}
            \item The header of the IP:
                  \begin{itemize}
                      \item version:\\
                            which version of IP it uses
                      \item IHL: \\
                            determines the length of the header
                      \item type of service: \\
                            indicates the service class
                      \item total length: \\
                            the joint length of the header and data part in bytes
                      \item identification: \\
                            every fragment of a datagram carries the same identification value.
                      \item DF: \\
                            ``don't fragment'' flag for the routers
                      \item MF: \\
                            ``more fragments'' flag has to be set in every fragment except the last.
                      \item fragment offset: \\
                            shows the place of the fragment within the datagram.
                      \item time to live: \\
                            its value should be decreased per second, at every hop they decrease its value by one
                      \item protocol: \\
                            contains the identifier of the transport-layer protocol
                      \item checksum: \\
                            a checksum on the header, used for handling errors produced by bad memory words inside the routers, which has to be recomputed at every hop
                      \item source address and destination address: \\
                            IP address
                      \item options: \\
                            left in for the extensibility of the next version.
                  \end{itemize}

                  \begin{figure}[H]
                      \centering
                      \includegraphics[width=0.8\textwidth]{../../18. Számítógépes hálózatok és Internet eszközök/img/ip_fejresz.png}
                      \caption{Header of IPv4}
                  \end{figure}

            \item IP address \\
                  Every host and every router on the Internet has an IP address that encodes the network number and the host number. The IP address is represented on 4 bytes. The IP is written in dotted decimal notation. (For example: 192.168.0.1) There are a few special addresses (Figure \ref{fig:ip_spec_cimek}).

                  \begin{figure}[H]
                      \centering
                      \includegraphics[width=0.8\textwidth]{../../18. Számítógépes hálózatok és Internet eszközök/img/ip_spec_cimek.png}
                      \caption{Special IP addresses}
                      \label{fig:ip_spec_cimek}
                  \end{figure}

                  Subnetworks:\\
                  The hosts in the same network have the same network number. A network can be divided into several parts for internal use, but to the outside world it appears as a single network. For identification the router needs to know the subnet mask. In the routing table at the routers there are entries of the form (network,0) and (own network, host). If there is no match, then the packet is forwarded toward the default router.\\

                  Depletion of IP addresses:\\
                  The IP addresses ran out fast. Solution: Classless Inter-Domain Routing (CIDR). The routing becomes more complex: Every entry is supplemented with a 32-bit mask. An entry can from now on be characterized by a triple: (ip address, subnet mask, output line). In case of a new packet the subnet address is masked out from the destination address, and in case of a match it is forwarded toward the longest match.

                  Another method is NAT, which is a quick fix for the problem of the depletion of IP addresses. For internet traffic every company is given one or at least few IP addresses, while within the company a unique IP address is used for every computer for internal routing: \\
                  10.0.0.0 - 10.255.255.255 : 16 777 216 unique addresses\\
                  172.16.0.0 - 172.31.255.255 : 1 048 576 unique addresses\\
                  192.168.0.0 - 192.168.255.255 : 65 536 unique addresses\\

                  IPv6:\\
                  In contrast to IPv4 it uses 16-byte addresses; we write it as 8 groups, each consisting of four-four hexadecimal digits. (For example: 8000:0000:0000:0000:0123:4567:89AB:CDEF) The IP header was simplified, which enables faster processing for the routers. A significant step was taken toward security.

                  \begin{figure}[H]
                      \centering
                      \includegraphics[width=0.8\textwidth]{../../18. Számítógépes hálózatok és Internet eszközök/img/ipv6.png}
                      \caption{Header of IPv6}
                  \end{figure}


        \end{itemize}
    \item[Protocols] \hfill
        \begin{itemize}
            \item Internet Control Message Protocol - ICMP \\
                  Its task is the reporting of unexpected events. Several kinds of ICMP message were defined:
                  \begin{itemize}
                      \item Destination unreachable
                      \item Time exceeded
                      \item Parameter problem
                      \item Source quench
                      \item Echo request
                      \item Echo reply
                      \item etc.
                  \end{itemize}
            \item Address Resolution Protocol - ARP \\
                  Its task is the matching of an IP address to a physical address. It sends out a ``Whose is the IP address 192.60.34.12?'' packet onto the Ethernet by broadcast on the subnetwork. Every host checks whether the IP address in question is theirs. If the IP coincides with the host's own IP, then it responds with its own Ethernet address.
            \item Reverse Address Resolution Protocol - RARP \\
                  Its task is the matching of a physical address to an IP address. The newly started station sends out a packet onto the Ethernet by broadcast: ``My 48-bit Ethernet address is 14.04.05.18.01.25. Does anyone know my IP address?'' And the RARP server responds with the appropriate IP address when it sees the request.
            \item Open Shortest Path First - OSPF \\
                  OSPF is responsible for routing within the ASes (Autonomous System). It maps the topology of the network, and senses the changes. It represents the topology with a weighted directed graph, in which it looks for cheapest paths.
            \item Border Gateway Protocol - BGP \\
                  Its task is that political considerations play a role in the routing decisions between ASes (e.g. Traffic starting or ending at IBM should not go through Microsoft, or Let us go through Albania only if there is no other path to the destination.)\\
                  From the point of view of a BGP router the world consists of ASes and the lines passing between them. (Two ASes are connected if there is a line connecting their BGP routers between them.)
                  From the point of view of transit traffic there are 3 kinds of network:
                  \begin{itemize}
                      \item Stub networks, which have only a single connection to the BGP graph
                      \item Multiply connected networks, which transit traffic could use, but they refuse this
                      \item Transit networks, which, with some restriction and generally for payment, are willing to handle third-party packets
                  \end{itemize}
        \end{itemize}
\end{description}
\section{Transport layer}
\begin{description}
    \item[Definition] \hfill \\
        The transport layer ensures that the data transfer between the users is transparent. The layer ensures, and checks, the reliability of a given connection. To the beginning of the data received from the application layer it attaches a so-called header, which indicates which transport-layer protocol the data is sent with. Some protocols are connection-oriented. This means that the layer tracks the data packets, and in case of an error takes care of the resending of the packet or packets.

    \item[Connectionless and connection-oriented] \hfill \\
        The connection-oriented protocol hides the possible errors of the transmission from the applications, we do not have to deal with the lost, or doubly arrived, or corrupted packets, nor with the order in which they arrived. However, this degrades its performance.

        In the connectionless case there is no need for breaking the data into frames, and there is no packet resending.
    \item[Reliability] \hfill \\
        The criteria of reliability:
        \begin{itemize}
            \item The arrival of every packet is acknowledged.
            \item The unacknowledged data packets are resent.
            \item A checksum is assigned to the header and the packet.
            \item It numbers the packets, and at the receiver the packets are sorted based on their sequence numbers.
            \item It deletes duplicates.
        \end{itemize}
    \item[Congestion control] \hfill \\
        On every network the transmission bandwidth is bounded. If we lead more data into the network, then it leads to congestion, or even to the collapse of the network (congestive collapse). Consequence: The data packets do not arrive.

        Avalanche phenomenon: \\
        The overloading of the network causes the loss of packets, which results in the resending of the packet. The resending further increases the load of the network so the packet loss will be even greater. This results in even more resent packets. ... \\
        The task of congestion control is the avoidance of the avalanche phenomenon.

        Requirements toward congestion control:
        \begin{itemize}
            \item Efficiency: \\
                  The transmission is large, while the delay is small.
            \item Fairness: \\
                  Every process gets an approximately equal share of the bandwidth. (The possibility of prioritization exists)
        \end{itemize}

        The tools of congestion control:
        \begin{itemize}
            \item Capacity increase
            \item Resource reservation and access control
            \item Load reduction and regulation
        \end{itemize}

        Strategies:
        \begin{itemize}
            \item Sliding window \\
                  Regulation of the data rate with the help of a window. The receiver determines the size of the window (wnd). If its receiving buffer is full, it decreases it to 0, otherwise it sends \textgreater0. The sender does not send more packets if the number of sent, not-yet-acknowledged packets reached the size of the window.
                  \begin{figure}[H]
                      \centering
                      \includegraphics[width=0.6\textwidth]{../../18. Számítógépes hálózatok és Internet eszközök/img/csuszoablak.png}
                      \caption{Sliding window}
                  \end{figure}
            \item Slow-start \\
                  The sender must not immediately use up the window size sent by the receiver. It determines another window (cwnd - Congestion Window), which it chooses. Then finally what it sends in: min\{wnd, cwnd\}. Initially cwnd = MSS (Maximum Segment Size). At every packet, after the received acknowledgment it increases: cwnd = cwnd + MSS (that is, it doubles after every RTT). This goes on until an acknowledgment is once missing.
            \item TCP-Nagle \\
                  It has to be ensured that the small packets get close to each other in time at delivery, and that in case of a lot of data the large packets get preference. \\
                  For this: The small packets do not get sent out while acknowledgments are missing (a packet is small if the data length \textless MSS). If the acknowledgment of the previously sent packet arrives, it sends the next one.
            \item TCP Tahoe and Reno
                  TCP uses both the sliding window and the Slow-start mechanism. Although the initial rate is small, the size of the window grows rapidly. When cwnd reaches the ssthresh (slow start threshold) value it switches into congestion avoidance state. TCP Tahoe and Reno are congestion avoidance algorithms. The two algorithms differ in how they detect and handle packet loss. \\
                  \underline{TCP Tahoe}: For the detection of congestion it sets a timer for the arrival of the expected acknowledgment.
                  \begin{itemize}
                      \item At connection setup: cwnd = MSS, ssthresh = $2^{16}$
                      \item At packet loss : Multiplicative decrease \\
                            cwnd = MSS, ssthresh = 	max\{2MSS, $\frac{min\{cwnd, wnd\}}{2}$\}
                      \item cwnd $\leq $ ssthresh : Slow-start \\
                            cwnd = cwnd + MSS
                      \item cwnd \textgreater ssthresh : Additive Increase \\
                            cwnd = 	cwnd + MSS$\frac{MSS}{cwnd}$
                  \end{itemize}
                  \underline{TCP Reno}: For the detection of congestion it uses both a timer and fast retransmit. [Fast retransmit: 3 acknowledgment duplicates arrive for the same packet (4 identical acknowledgments), then it resends the packet and enters the Slow-start phase.] \\
                  After fast retransmit: ssthresh = max\{$\frac{min\{wnd,cwnd\}}{2}$, 2MSS\}, cwnd = ssthresh + 3MSS. \\
                  Fast recovery after fast retransmit increases the rate after every further acknowledgment: cwnd = cwnd + MSS.
        \end{itemize}

        Efficiency and Fairness: \\
        The transmission is maximal if the load almost reaches the capacity of the network. If the load grows further, the buffers overflow, packets get lost, have to be resent, the response time drastically increases. This is called congestion. Therefore instead of the maximum load, it is advisable to set the load of the network near the knee. Here the response time rises only slowly, while the data transmission is already near the maximum

        A good congestion avoidance strategy keeps the load of the network near the knee: \textit{efficiency}. In addition it is important that we serve every participant with an equal rate: \textit{fairness}

        Let the data rate of the $i$-th participant at time $t$ be denoted by $x_i(t)$.
        Every participant updates its data rate in the $(t+1)$-th round:
        \begin{align*}
            x_i(t+1) = f_0(t) \quad  \text{if} \ \sum_{i=1}^{n}x_i(t) \leq K \\
            x_i(t+1) = f_1(t) \quad  \text{if} \ \sum_{i=1}^{n}x_i(t) > K
        \end{align*}
        where $f_0(x) = a_I + b_Ix$ is the increase, $f_1(x) = a_D + b_Dx$ the decrease strategy.

        Special cases:
        \begin{itemize}
            \item \textbf{M}ultiplicative \textbf{I}ncrease \textbf{M}ultiplicative \textbf{D}ecrease - MIMD:
                  \begin{align*}
                      f_0(x) = b_Ix \quad (b_I > 1) \\
                      f_1(x) = b_Dx \quad (b_D < 1)
                  \end{align*}
            \item \textbf{A}dditive \textbf{I}ncrease \textbf{A}dditive \textbf{D}ecrease - AIAD:
                  \begin{align*}
                      f_0(x) = a_I+x \quad (a_I > 0) \\
                      f_1(x) = a_D+x \quad (a_D < 0)
                  \end{align*}
            \item \textbf{A}dditive \textbf{I}ncrease \textbf{M}ultiplicative \textbf{D}ecrease - AIMD:
                  \begin{align*}
                      f_0(x) = a_I+x \quad (a_I > 0) \\
                      f_1(x) = b_Dx \quad (b_D < 1)
                  \end{align*}
        \end{itemize}

    \item[Multiplexing, demultiplexing] \hfill \\
        By multiplexing in telecommunications is meant the procedure when two or more channels are gathered into one channel in such a way that with the inverse multiplexing operation, or demultiplexing, or demuxing, the original channels can be produced. The original channels can be identified with a so-called coding scheme.
    \item[Interaction models] \hfill
        \begin{itemize}
            \item Bidirectional byte stream \\
                  The data is transmitted as two opposite-direction byte sequences. The content is not interpreted. The temporal behavior of the data packets can change: the transmission speed can increase, decrease, with different delay, they can also arrive in a different order. It tries to deliver the data packets close to each other in time. It tries to use the transmission medium efficiently.

            \item RPC \\
                  For accessing a procedure to be run on a remote machine network communication is needed, this procedure-call mechanism is hidden by RPC (Remote Procedure Call).

                  The steps of the call:
                  \begin{enumerate}
                      \item The client process locally calls the
                            client stub.
                      \item It packs the identifier and parameters of the procedure,
                            calls the OS.
                      \item It sends the message over to the remote OS.
                      \item It passes the message to the
                            server stub.
                      \item It unpacks the parameters,
                            passes them to the server.
                      \item The server locally calls the procedure,
                            gets the return value.
                      \item The sending back of this to the client
                            happens similarly, in the reverse direction.
                  \end{enumerate}
                  \begin{figure}[H]
                      \centering
                      \includegraphics[width=0.6\textwidth]{../../18. Számítógépes hálózatok és Internet eszközök/img/rpc.png}
                      \caption{RPC}
                  \end{figure}
        \end{itemize}
    \item[Protocols] \hfill
        \begin{itemize}
            \item TCP
                  \begin{itemize}
                      \item Creating a reliable data stream between two endpoints
                      \item It divides the application layer's data flow into packets
                      \item The other endpoint sends an acknowledgment of the reception of the packets
                  \end{itemize}
                  The content of the TCP header:
                  \begin{itemize}
                      \item source port (16 bit)\\
                            Identifies the sender process
                      \item destination port (16 bit)\\
                            The identifier of the addressee process
                      \item sequence number (32 bit)\\
                            The sequence number of the first data byte within the current segment. If the SYN flag bit has value 1, then this sequence number is the initial sequence number, that is, the sequence number of the first data byte will be the initial sequence number + 1.
                      \item acknowledgment number (32 bit)\\
                            If the ACK flag bit has value 1, then it contains the sequence number that the receiver wishes to receive next. It is sent at every connection setup.
                      \item header length (4 bit)\\
                            The length of the TCP header in 32-bit units.
                      \item Window (16 bit)\\
                            How many bytes can be sent starting from the acknowledged byte. (The value 0 is also valid.)
                      \item Checksum (16 bit)\\
                            For checking the data, header, and pseudo-header part.
                      \item Options (0-40 byte)\\
                            Designed for possibilities outside the standard header. The most important such possibility is the specification of the MSS, that is, the maximum segment size. Further options: MD5 signature, TCP-AO, ``usertimeout'', etc.
                      \item urgent pointer (16 bit) \\
                            Indicates the place of the urgent data measured in bytes relative to the current byte sequence number.
                      \item Flag bits (6)
                            \begin{enumerate}
                                \item URG -- Urgent flag bit.
                                \item ACK -- acknowledgment indication.
                                \item PSH -- Indicates that fast data forwarding is needed for the user layer.
                                \item RST -- Indicates the one-sided breaking of the connection.
                                \item SYN -- Indicates sequence number synchronization.
                                \item FIN -- Indicates the end of the data stream.
                            \end{enumerate}
                  \end{itemize}
                  \begin{figure}[H]
                      \centering
                      \includegraphics[width=0.8\textwidth]{../../18. Számítógépes hálózatok és Internet eszközök/img/tcp_fejlec.png}
                      \caption{TCP Header}
                  \end{figure}
                  Characteristics of TCP:
                  \begin{itemize}
                      \item Connection-oriented
                      \item Reliable
                      \item Bidirectional byte stream
                  \end{itemize}
            \item UDP
                  \begin{itemize}
                      \item A simple, unreliable service for sending packets
                      \item The application layer determines the size of the packet
                      \item It transforms the input into a datagram
                  \end{itemize}
                  A connectionless protocol. It uses for the transmission segments that consist of an 8-byte header and the user data.\\
                  The header contains:
                  \begin{itemize}
                      \item a source port (2 byte);
                      \item a destination port (2 byte);
                      \item a UDP segment length value (2 byte);
                      \item a UDP checksum (2 byte)
                  \end{itemize}
                  UDP does not perform flow control, error handling or resending after the reception of a bad segment. In the case of client-server applications UDP can be especially useful because of the short messages.
        \end{itemize}
\end{description}

\section{Application layer}
\begin{description}
    \item[DNS] \hfill \\
        The Domain Name System (DNS) is a hierarchical, highly distributed naming system for computers, services, and any resource connected to the internet or a private network. To the domain names assigned to the participating entities it associates various information. As its most important function, it ``translates'', ``resolves'' the domain names meaningful to humans into numeric identifiers understandable to network devices, with the help of which these devices can be found, addressed on the network.
        A frequently used analogy for explaining the domain name system is that it is a kind of telephone book of the internet, from which one can look up the IP addresses belonging to the human-interpretable computer host names. For example, to the domain name www.example.com belong the addresses 192.0.32.10 (IPv4) and 2620:0:2d0:200::10 (IPv6).

        DNS makes it possible to assign names to groups of internet resources in a way that does not depend on the physical place of the resources. Thus the World Wide Web (WWW) hyperlinks, internet contact data can stay consistent and constant even if a change occurs in the internet's routing system, or the participant uses a mobile device. A further purpose of internet domain names is simplification, a domain name (e.g. www.example.com) is much easier to remember than an IP address such as 208.77.188.166 (IPv4) or 2001:db8:1f70::999:de8:7648:6e8 (IPv6). The users can thus remember the web (URL) and e-mail addresses meaningful to them, without knowing how the computer actually reaches these.

        In DNS the responsibility for the assignment of domain names and the assignment of IP addresses is delegated; to every domain belongs an authoritative name server. The authoritative name servers are responsible for their own domains. This responsibility can be further delegated, so for the sub-domains another name server can be responsible. This mechanism is behind the distributed and fault-tolerant operation of DNS, and therefore it is not necessary to maintain and constantly update a single central directory.

        In the domain name system other information is also stored, for example a list of the mail servers receiving e-mail for a given internet domain. As a worldwide, distributed, keyword-based redirection service, the Domain Name System is a fundamentally important element of the functionality of the internet.

        The DNS database can also be used for storing RFID tags, UPCs, IP telephone numbers and many other things.

        The Domain Name System specifies the technical capabilities of the database, and in addition describes the DNS protocol forming part of the internet protocol family, specifies in detail the data structures and communication used in DNS.

    \item[HTTP] \hfill \\
        HTTP (HyperText Transfer Protocol) is an information transfer protocol for distributed, collaborative, hypermedia, information systems.

        The development of HTTP was coordinated by the World Wide Web Consortium and the Internet Engineering Task Force in the form of RFCs. RFC 2616, issued in 1999, defines HTTP/1.1, which was replaced by the end of 2015 by the HTTP/2.0 version, which is defined by RFC 7540. Officially this is the latest protocol.

        HTTP is a request-response based protocol between client and server. The HTTP clients are also usually referred to by the collective name ``user agent''. The user agent is typically, but not necessarily, a web browser.

        HTTP is located above the TCP/IP layer. HTTP can also be implemented over another reliable transport layer, whether on the internet or another network. It exclusively uses the TCP protocol, since data loss is not permissible.
    \item[DHCP] \hfill \\
        The Dynamic Host Configuration Protocol (DHCP) is a computer network communication protocol.

        This protocol solves the problem that the network endpoints (for example computers) connecting to the TCP/IP network automatically get the settings needed for using the network. Such usually are, for example, the IP address, network mask, default gateway, etc.

        DHCP is a server-client based protocol, broadly it consists of the DHCP requests sent by the clients, and the DHCP responses given by the server.

        With DHCP, IP addresses can be assigned dynamically, so the IP addresses of computers disconnecting from the network are received by the computers connecting to the network, therefore the narrower address ranges can be used more efficiently.

        3 kinds of IP assignment are possible with DHCP:
        \begin{itemize}
            \item manual (based on MAC address)
            \item automatic (by specifying the IP range that can be issued with DHCP)
            \item dynamic (by specifying an IP range, but with the ``reuse'' of the IP addresses)
        \end{itemize}

    \item[ARP] \hfill \\
        ARP (Address Resolution Protocol) is, in computer science, a method used on computer networks for matching IP addresses and physical addresses to each other. In practice, knowing the IP address we get the 48-bit physical address determined by the manufacturer of the network card. Because of the widespread nature of IPv4 and Ethernet it is generally used for translation between IP addresses and Ethernet addresses, but it is also operational in ATM or FDDI networks.

        Two client machines use the ARP protocol in the following four basic cases:
        \begin{itemize}
            \item If the two client machines are on the same network, and one wants to send a packet to the other.
            \item If the two client machines are on different networks, and reach each other through a gateway/router.
            \item If a router has to forward a client's packet through another router.
            \item If a router has to forward a client's packet to the addressee, which is on the same network.
        \end{itemize}

        In the first case the two clients are on the same physical network (so they can communicate directly with each other without using a router). The other three cases are the most common on the internet, where generally any two computers are more than 3 hops away from each other.

        Let us imagine that computer A sends a packet to computer D, and between them are routers B and C. In the 2nd case A sends to B, in the 3rd case B sends to C, and in the 4th case C sends a packet to D.
\end{description}

\end{document}
